\subsection{Questão 59}

\subsubsection{Enunciado}

Na Fig.~\ref{fig:33-59} um feixe luminoso que se propaga inicialmente no material 1 é refratado para o material 2, atravessa esse material e incide com o ângulo crítico na interface dos meios 2 e 3.

Os índices de refração são
\[
n_1 = 1{,}60,
\qquad
n_2 = 1{,}40
\qquad \text{e} \qquad
n_3 = 1{,}20.
\]

\begin{enumerate}
    \item[(a)] Qual é o valor do ângulo \(\theta\)?

    \item[(b)] Se o valor de \(\theta\) aumenta, a luz consegue penetrar no meio 3?
\end{enumerate}

\begin{figure}[h]
    \centering
    \imagemquestao{figures/fig33-59.jpeg}
    \label{fig:33-59}
\end{figure}


\subsubsection{Resolução}


\vspace{1cm}
